\section{Hercules’ Ninth – The Girdle of Hippolyta}

Thucydides the Athenian notes that of pre-history, he is obliged to accept what the poets say, although certainly this is unsatisfactory, and the poets are not to be trusted. Since Thucydides was an Athenian, he presumably shared in the general blase attitude which they took towards the rites and rituals. The Spartans would postpone invasions if the auguries were inauspicious, whereas the Athenians did not hesitate to tear down temples dedicated to the gods and fortify them as watch towers. The Spartans consulted the Delphic oracle before waging war, and the Athenians held a democratic (or oligarchic) debate. So we see that the rationalizing outlook is nothing new, for the Athenians defended themselves from charges of impiety by using the excuse that they were compelled by war's necessity, which was a general law of nature, and that the gods would understand this. They didn't reject the gods, they merely re-interpreted them, since occasionally a philosopher could offend even their sensibility enough to run the risk of impiety charges. In fact, with their maritime Empire, democracy, riches, and political expediencies (see the Melian dialogue\footnote{\url{https://en.wikipedia.org/wiki/Melian\_dialogue}}), the Athenians seem to be a very ``modern'' people indeed.

Those who struggle to undertake Hercules' labors should take heart that these trends and currents against which they swim are in no way a novelty. Did I mention that Athens lost the Peloponnesian War against Sparta? It was the hubris of their Sicilian Expedition which did them in. Hubris is always present when the gods are disrespected. \emph{Those whom the gods wish to destroy, they first make mad}. Immersing yourself in Tradition is one way of ``not going mad''; the ancients knew that poetry, myth, legend, and religion were important ways of practicing waking up before one was even fully awake. At least if one goes ``mad'' during the epicycle of Tradition, there is a cure available within the cycle, since it is ``all of a piece''. As we see with Hercules, making human detours or mortal mistakes do not derail him on his path. This is not merely because God ``favors'' him, but because Hercules is acting a part -- that is, he is ``acting'' above himself, rather than mimicking that which is below him. He is rising to his true Self, doing the opposite of what a play actor does. If you tell a man, ``God is Light'', that still requires either a) an epiphany on his part, or b) more explanation. But if you tell him, that in order to go to heaven, he has to pass the three headed dog and travel by the path that borders hell, in the shadow of the dark forest, and save the maiden who waits for him, then he can begin to understand.

In this episode, our hero does quite a bit of killing. As I've tried to explain to some liberal friends, if it's worth dying for, it's probably worth killing for in some way, as well. Although I don't fully understand what Christ meant by ``those who live by the sword, die by the sword'', there are other Bible proof texts, as long as people are asking.

\begin{quotex}
Do not think that I have come to bring peace to the earth; I have not come to bring peace but a sword\footnote{\url{http://en.wikipedia.org/wiki/But\_to\_bring\_a\_sword}}.
\flright{\textsc{Matthew 10:34}\footnote{\url{http://bibref.hebtools.com/?book=\%20Matthew&verse=10:34&src=\%21}}}

I came to bring fire to the earth and how I wish it were already kindled! Do you think that I have come to bring peace to the earth? No, I tell you, but rather division.
\flright{\textsc{Luke 12:49-51}\footnote{\url{http://bibref.hebtools.com/?book=\%20Luke&verse=12:49-51&src=\%21}}}

And he said unto them, When I sent you without purse, and scrip, and shoes, lacked ye any thing? And they said, Nothing. Then said he unto them, But now, he that hath a purse, let him take it, and likewise his scrip: and he that hath no sword, let him sell his garment, and buy one\footnote{\url{http://en.wikipedia.org/wiki/Sell\_your\_cloak\_and\_buy\_a\_sword}}. For I say unto you, that this that is written must yet be accomplished in me, And he was reckoned among the transgressors: for the things concerning me have an end. And they said, Lord, behold, here are two swords. And he said unto them, It is enough.
\flright{\textsc{Luke 22:35-38}\footnote{\url{http://bibref.hebtools.com/?book=\%20Luke&verse=22:35-38&src=\%21}}}
\end{quotex}


Hercules will end up adventuring among the Amazons; Hera comes down, and stirs up their ire, as it would appear that Hercules' power and charm (by itself) would have overcome them. Only Hippolyta is untouched, but she is doomed in her act of going to Hercules to aid him. Did Hera's rage affect our hero, as well? We think that a touch of it did cloud him, here.

\begin{quotex}
Diodorus Siculus\footnote{\url{http://en.wikipedia.org/wiki/Diodorus\_Siculus}}\footnote{\url{http://en.wikipedia.org/wiki/Amazons\#cite\_note-29}} enlists nine Amazons who challenged Heracles to single combat during his quest for Hippolyta's girdle and died against him one by one: Aella\footnote{\url{http://en.wikipedia.org/wiki/Aella\_\%28Amazon\%29}}, Philippis, Prothoe, Eriboea, Celaeno\footnote{\url{http://en.wikipedia.org/wiki/Celaeno}}, Eurybia, Phoebe\footnote{\url{http://en.wikipedia.org/wiki/Phoebe\_\%28mythology\%29}}, Deianeira, Asteria\footnote{\url{http://en.wikipedia.org/wiki/Asteria}}, Marpe, Tecmessa, Alcippe. After Alcippe's death, a group attack followed.
\end{quotex}


In another version, Hercules' takes Melinippe (Hippolyta's sister) captive, and successfully exchanges the girdle in return for her life.

For most versions, after the mob attack on Hercules (which fails), Hippolyta storms to his rescue, but Hercules thinks she has betrayed him, and strikes at her with the inevitable result.

The Amazons are connected with the city of Troy, \& came to its defense during the siege. Coupling this with the episode of Dido\footnote{\url{http://en.wikipedia.org/wiki/Dido\_\%28Queen\_of\_Carthage\%29}} \&{} Aeneas (in which Aeneas inadvertently kills Dido and her love by leaving), we can say that the Amazon theme has a strong association with the trajectory and rise of the hero. They come after the affliction, \& before the exultation and triumph. The rise of Roma is under the occult name of Love.

Hera is clearly a retarded element in the myth, but the Amazons are not. The girdle had been given by Ares to the queen, and represents feminine power to stand as a man against men. When Hippolyta exchanges the girdle for her sister, she is returned to the world of sisterly and feminine existence. She is tempered and made flexible. Hercules' death blow may or may not be literal: there is an exchange between them which destroys Hippolyta's warlike status \& leads to Hercules' obtaining of the girdle.

We have to note that Hercules is not seeking the girdle out of hatred for the Amazons. It is part of the given quest, \& Hercules isn't going to boast about it or even wear it. Hercules' adventures were the subject of a good deal of contemplation during the Middle Ages\footnote{\url{http://books.google.com/books?id=H5g08PMrkjgC&pg=PT161&lpg=PT161&dq=hercules+labors+christ&source=bl&ots=pgNOIWGurz&sig=W39hmhIerqtKjEUo4W4eY763-Xk&hl=en&sa=X&ei=OGF8UtWoMOzIsASgvYLgBw&ved=0CFwQ6AEwCQ\#v=onepage&q=hercules\%20labors\%20christ&f=false}}. The allegorization of Hercules was not restricted to Christian thinkers, \& there were even Tarot cards which used the Hercules motif\footnote{\url{http://pre-gebelin.blogspot.com/2008/04/prodicus-allegory-of-virtue-and-vice.html}}. In general, we may fairly say that Hercules has by right always been understood as symbolizing the strong man of the real self, who overcomes the lower self and even its own self. Indeed, the theme remained popular through the Baroque era, up until very modern times\footnote{\url{http://www.ourcivilisation.com/smartboard/shop/fowlerjh/chap1.htm}}.

Why make much of the ``true Self'' and ``false Self''? Well, for one thing, we don't see Hercules sleeping his way through the adventures by bedding a bunch of females. At the most, we reach any hint of that in this story, and only a hint. We don't find him killing indiscriminately or widely (during this adventure, he does kill some men, but desists when kinsmen of the king are offered in exchange for his fallen comrades, and the Amazon battle is self-preservation). He is not ``tempted'' in any conventional way, with delights or ease. So the medieval and pre-medieval interpretation is legitimate: Hercules is not a masterless man.

Alice Bailey drew parallels\footnote{\url{http://www.souledout.org/hercules/herculesvirgo/hercvirgo.html}} with the Zodiac \& the meaning of Virgo, one of the oldest signs. What comes after his death-blow to Hippolyta is instructive:

\begin{quotex}
Hippolyte is the Queen of the Amazons, a great race of women warriors, to whom Venus, the Queen of Love, has given a girdle, the symbol of unity, motherhood and the sacred Child, achieved through struggle. Traveling to the shores of her kingdom, Hercules fights Hippolyte and kills her for the girdle. Too late, he realizes that Hippolyte had not put up any struggle but offered the girdle to Hercules freely, under instruction from God. In a state of shock, Hercules sets forth to redeem this failure. While journeying back to the shores of the sea, he hears the screams of the beautiful Hesione being attacked and then swallowed by a sea monster. Hercules swims furiously out to sea to rescue her and, as the monster turns on him, opening its mouth to roar, Hercules rushes down the red tunnel of its throat, grasps Hesione and hews his way out of the serpent's belly. Thus, after slaughtering the feminine principle that gave him what he needed, Hercules makes amends by rescuing that which also needed him.
\end{quotex}


What is being uncovered here is the falseness of the corrupt feminine principle (Hera and the mob Amazons) and the new Woman called into being by the progress of the Hero: the man is acquiring firmness without brutality, and the woman is acquiring docility and flexibility without dishonor. Man becomes divine, Woman becomes eternal, together, but in a different way.

Hippolyta is a perfect Dido, and is moreover the perfect woman (although belatedly). Hercules responds by essentially becoming a Christian knight. With this adventure, we have entered the world of the Middle Ages and their romances, deliberately satirized in \emph{Don Quixote}. But rather than becoming Orlando Furioso\footnote{\url{http://en.wikipedia.org/wiki/Orlando\_Furioso}}, Hercules acts directly. In literary words, Hercules takes the path of Dante rather than the path of Faust; he is open to vertical change, and avoids meandering and wallowing in self-pity, regret, manipulation of Nature, or stasis.

In the liberal world, we are offered a ``God Without Thunder\footnote{\url{http://spider.georgetowncollege.edu/htallant/border/bs7/kennedy.htm}}``. In this world, God must act according to our preconceptions. Thus, YHWH is guilty of genocide and the worst forms of ancient prejudice and violence, and is a veritable demon, whereas Jesus becomes a prophet-sage who aimed at freeing man from all darkness with the news that ``God is Light''. This is the new civil religion for our time. It delights in opposing false anti-theses, in order to develop facetious syntheses. The opposite of this is to peer deeper into the mysteries of existence, \& this is done, not through rationalization, but (for most people) legend, myth, history, poetry, and religion. These traditional forms are analogue and temporary substitutes for Right or Pure Reason: they are ``practice'', like children playing at sword-sticks -- it is also deadly serious, but a form of play. The first card of the Tarot enjoins us to turn ``work into play''.

For Hercules the warrior knows there are monsters at large, not least of all in his own heart, \& that they must all be conquered, slain, or transformed, as may be. Evil exists, \& there is really no way to get God off the hook, rationally speaking. Either one begins the path by piety, submission, and respect for the higher power which tasks you (in which case these higher powers differentiate themselves into Hera/adversary and Zeus/protector, and the end of which is apotheosis and divinization, in which the hero is united to the separated divine) or else you try clever and abstract mental gymnastics. The liberal worldview has chosen the latter, \& it is far easier to simply rewrite, ignore, or ridicule Tradition than to take up the cross, and follow the path of the hero-turning-knight. Is there a third way? I do not believe that there is. You can only see more deeply into what is already true, what is already old, what is already there. There is no revolutionary resolution or higher truth that preserves the profundity of Tradition, while somehow transcending the thunder of the Gods which offend our sensibilities. One will devour the other, and this itself shocks the sensibilities!

The temptation here is to mental passion and pride, which is why Gornahoor for Lent last year recommended a ``mind fast'' rather than a food fast. The myth clearly walks the initiate through the cycle of the Zodiac, and includes not only detours and failures, but also conspiracies and meta-conspiracies, as well as a varied array of monsters. It is, in other words, psychologically and spiritually both replete and complete. It can stand alone as its own veritable natural Gospel, were the coming of Christ hundreds of years off, because it is the primordial Truth in the first place.

When Hercules walks away with the girdle, but then goes and saves the maiden on his way home, \& then surrenders the war-relic to his adversary as proof of passing the tests, he has ``learned justice and mercy'' by triumphing over the passions of vengeance and lust and the temptation to use the rod of iron in the cause of human wrath or personal gain. In each labor, he proves that he has deeply assimilated the implications of each previous lesson learned, even as he stumbles.

The point is not to avoid failure, but to get back up. That was always the point to begin with. The quest doesn't end, until Hercules is done acting on each lesson learned. Fate is now weaving itself around the hero, rather than weaving the hero around Itself. \emph{Ducunt volentem Fata, nolentem trahunt}. Hercules, exposed warts and all\footnote{\url{http://www.theguardian.com/artanddesign/2013/nov/08/cromwell-portraitist-samuel-cooper-exhibition}}, is now ta'veren\footnote{\url{http://wot.wikia.com/wiki/Ta\%27veren}}. His deeds are increasingly judged only by his own conscience, \& by God. The shallow and base will call out (as they do against God) ``you have waded through slaughter to a throne!''. But Zeus \& Hercules \& Hippolyta know that he is becoming a parfait and gentle\footnote{\url{http://www.luminarium.org/medlit/knightport.htm}} knight.

What does it mean? As late as the last adventure, we saw the growing conspiracy, which comes out into the open now, using an arm of humanity (the mob feminists of the day) against the hero in a berserk manner. Hercules stays focused on learning personal lessons, and will not allow himself to be cowed or bullied either inside or out. He stays focused on the lesson, even when failing it. And because of that, he cannot fail.

Is this not the hidden meaning of ``all things work to the good of them who love the Lord''?


\flrightit{Posted on 2013-11-15 by Logres}

\begin{center}* * *\end{center}

\begin{footnotesize}
\begin{sffamily}
\texttt{IA on 2013-11-18 at 08:59 said:}

Hi,

As Ezra Pound said, culture is what's left over after you forgot what you tried to learn. Culture is instinct but learned instinct. Art and myth reconnect us to something we can only glimpse in the corner of our eye. When we look directly at it, it disappears. Humans need mystery, evidently.

Religion and high culture satisfy this craving for mystery. And, let's face it, the world is a pretty strange place. But, in an impious age images and gestures can only signify emptiness and irony, with no storytelling. High preists of modernism mystify and baffle with abstraction, discontinuity and strange words that allow secret knowledge to initiates. I think its gnosticism.

Anyway, thanks.


\hfill

\end{sffamily}
\end{footnotesize}

